\Needspace{14\baselineskip}
\section{Responses to the comments of reviewer 1}
\label{sec:responses-reviewer1}
\begin{enumerate}
  \Needspace{5\baselineskip}
  \item \reviewercomment{Define N in the averaging equation and state which readings are included.}
    \response{The Method section now defines N as the positive number of readings in one batch and states that these readings measure the same quantity in the same units.}
  \Needspace{5\baselineskip}
  \item \reviewercomment{Give the measurement units and show the calculation for readings 10, 12, and 14.}
    \response{The Results section now identifies temperature in degrees Celsius and shows the calculation (10+12+14)/3=12 degrees Celsius.}
  \Needspace{5\baselineskip}
  \item \reviewercomment{Discuss the effect of an outlier and state the limits of the method.}
    \response{The new Limitations section states that an outlier can shift the mean and explicitly avoids claiming robustness to outliers.}
\end{enumerate}
